% ===== PRÉAMBULE CANONIQUE — à coller VERBATIM en tête de CHAQUE .tex =====
\documentclass[11pt,a4paper]{article}
\usepackage[utf8]{inputenc}\usepackage[T1]{fontenc}\usepackage{lmodern}
\usepackage[french]{babel}
\usepackage{amsmath,amssymb,mathtools}\usepackage{icomma}
\usepackage{siunitx}\sisetup{locale=FR}  % \qty{}{}, \si{}, \ang{}
\usepackage{xcolor}\usepackage{colortbl}
\usepackage{tikz}\usepackage{pgfplots}\pgfplotsset{compat=1.18}
\usepackage[siunitx,europeanresistors]{circuitikz} % circuits (élec.)
\usepackage[most]{tcolorbox}\usepackage{enumitem}\usepackage{array,booktabs}
\usepackage[a4paper,margin=2.2cm]{geometry}
\usetikzlibrary{arrows.meta,calc,angles,quotes,positioning,patterns,%
  backgrounds,shapes.geometric,decorations.pathmorphing,decorations.markings,intersections}
\definecolor{primary}{RGB}{222,49,99}\definecolor{primarydark}{RGB}{150,22,55}
\definecolor{defcol}{RGB}{0,114,178}\definecolor{methcol}{RGB}{52,92,148}
\definecolor{excol}{RGB}{0,135,90}\definecolor{attncol}{RGB}{200,40,55}
\tcbset{boxbase/.style={enhanced,breakable,boxrule=0.9pt,arc=2.5pt,
  left=3mm,right=3mm,top=2.3mm,bottom=2.3mm,fonttitle=\bfseries,coltitle=white}}
% --- boîtes TUTORIEL (noms EXACTS, ne pas renommer) ---
\newtcolorbox{cle}[1][]{boxbase,colback=defcol!6,colframe=defcol,
  title={\textsc{À retenir}\ifx\relax#1\relax\relax\else\textmd{ : \itshape #1}\fi}}
\newtcolorbox{methode}[1][]{boxbase,colback=methcol!7,colframe=methcol,sharp corners,
  title={\textsc{Méthode}\ifx\relax#1\relax\relax\else\textmd{ --- \itshape #1}\fi}}
\newtcolorbox{exemplebox}[1][]{boxbase,colback=excol!5,colframe=excol,
  title={\textsc{Exemple}\ifx\relax#1\relax\relax\else\textmd{ : \itshape #1}\fi}}
\newtcolorbox{piege}[1][]{boxbase,colback=attncol!6,colframe=attncol,
  title={\textsc{Piège}\ifx\relax#1\relax\relax\else\textmd{ : \itshape #1}\fi}}
\newtcolorbox{remarque}[1][]{boxbase,colback=black!4,colframe=black!45,
  title={\textsc{Remarque}\ifx\relax#1\relax\relax\else\textmd{ : \itshape #1}\fi}}
% --- boîtes EXERCICE / CORRIGÉ (noms EXACTS, auto-comptées) ---
\newtcolorbox[auto counter]{exo}[1][]{boxbase,colback=primary!4,colframe=primary,
  title={\textsc{Exercice}~\thetcbcounter\ifx\relax#1\relax\relax\else\textmd{ --- \itshape #1}\fi}}
\newtcolorbox[auto counter]{corrige}[1][]{boxbase,colback=excol!4,colframe=excol,
  title={\textsc{Corrigé}~\thetcbcounter\ifx\relax#1\relax\relax\else\textmd{ --- \itshape #1}\fi}}
\newcommand{\vv}[1]{\vec{#1}}\newcommand{\ds}{\displaystyle}
\newcommand{\R}{\mathbb{R}}\newcommand{\N}{\mathbb{N}}\newcommand{\Z}{\mathbb{Z}}
% =========================================================================
\begin{document}
\sisetup{per-mode=symbol}

\begin{exo}[la force de réaction]
Un corps $A$ exerce sur un corps $B$ une force $\vv F_{A\to B}$ d'intensité \qty{15}{\newton}.
\begin{enumerate}[label=\alph*)]
  \item Quelle est l'intensité de la force $\vv F_{B\to A}$ exercée par $B$ sur $A$ ?
  \item Précise sa direction et son sens par rapport à $\vv F_{A\to B}$, et sur quel corps elle s'applique.
\end{enumerate}
\end{exo}

\begin{exo}[deux patineurs se repoussent]
Deux patineurs se repoussent : chacun subit la même force $F=\qty{60}{\newton}$. Le patineur $A$ a une
masse $m_A=\qty{60}{\kilogram}$, le patineur $B$ une masse $m_B=\qty{40}{\kilogram}$.
\begin{enumerate}[label=\alph*)]
  \item Calcule l'accélération de $A$.
  \item Calcule l'accélération de $B$.
\end{enumerate}
\end{exo}

\begin{exo}[rapport des accélérations]
Un canon de masse $m_{\text{canon}}=\qty{1000}{\kilogram}$ tire un boulet de
$m_{\text{boulet}}=\qty{10}{\kilogram}$. La force exercée sur le boulet est, en réaction, la même que
celle exercée sur le canon.
\begin{enumerate}[label=\alph*)]
  \item Écris le rapport $\dfrac{a_{\text{boulet}}}{a_{\text{canon}}}$ en fonction des masses.
  \item Calcule sa valeur numérique.
\end{enumerate}
\end{exo}

\begin{exo}[qui accélère le plus ?]
Un poids lourd et une petite voiture entrent en collision : ils exercent l'un sur l'autre une force de
même intensité.
\begin{enumerate}[label=\alph*)]
  \item Lequel des deux subit la plus grande accélération ? Justifie par $a=F/m$.
  \item Les forces subies par les deux véhicules ont-elles la même intensité ?
\end{enumerate}
\end{exo}

\begin{exo}[action et réaction dans la vie courante]
Pour chaque situation, indique la force d'\emph{action} et sa \emph{réaction}.
\begin{enumerate}[label=\alph*)]
  \item Un marcheur : son pied pousse le sol vers l'arrière.
  \item Un nageur : ses bras poussent l'eau vers l'arrière.
\end{enumerate}
\end{exo}
\end{document}
